\documentclass[11pt]{amsart}

\usepackage[T1]{fontenc}
\usepackage[utf8]{inputenc}
\usepackage{lmodern}
\usepackage{amssymb}
\usepackage{mathtools}
\usepackage{booktabs}
\usepackage[margin=1in]{geometry}
\usepackage[hidelinks]{hyperref}

\setlength{\emergencystretch}{3em}

\newcommand{\lean}[1]{\texttt{\detokenize{#1}}}
\newcommand{\fn}[1]{\mathrm{#1}}
\newcommand{\leanref}[1]{\par\nopagebreak\smallskip
  {\small\raggedright\noindent Lean: #1\par}\medskip}

\newtheorem{theorem}{Theorem}
\newtheorem{lemma}[theorem]{Lemma}
\theoremstyle{definition}
\newtheorem{definition}{Definition}

\title[A verified jump gear above $2^{64}$]{A formally verified jump gear for segmented
sieving above $2^{64}$}
\author{Artur Flamandzki}
\date{}

\begin{document}

\begin{abstract}
A segmented sieve of Eratosthenes on a window $[F, F+W)$ far above $2^{64}$ spends most of its
time on sieving primes larger than the window, which hit it at most once and usually not at
all. The \emph{jump gear} of the PrimeAtlas engines v7 handles these numbers without computing
$F \bmod q$ for each of them: on blocks of $q$ it replaces $F/q$ by a tangent line, turns the
landing condition into a rotation on $\mathbb{Z}/2^{64}$ in 64.64 fixed point, finds the first
candidate by a Euclid-like recursion and steps through the others by the three-gap theorem.
The mathematical core of the method is Helfgott's improved sieve of
Eratosthenes~\cite{helfgott}. We prove in Lean~4 / Mathlib that the C code of the gear, as
written, visits every landing of a block and computes its offset exactly, for engines v7.0,
v7.1 and v7.2. The model covers 128- and 64-bit wraparound, the truncations of the
fixed-point quantities, the tangent error, and a threshold computed in IEEE-754 double
precision under the standard rounding model. The proof gives the engine bound
$30F < 2^{128}$, and a concrete instance, checked by the kernel, shows that the engines'
earlier guard $F < 2^{126}$ let a 128-bit product overflow; the engines now use the proved
bound. A certificate proved along the way lets a block keep the three-gap steps of the
previous one; engine v7.4 uses it. We also prove that the engines' gear threshold and block
factor are the optimum of their cost model. Every statement in the results section is a
theorem of the accompanying Lean formalization, which builds with no \lean{sorry} and uses
only the axioms \lean{propext}, \lean{Quot.sound} and \lean{Classical.choice}.
\end{abstract}

\maketitle

\section{Introduction}

A segmented sieve of Eratosthenes~\cite{bayshudson} marks, in a window $[F, F+W)$, the
multiples of every prime up to $\sqrt{F+W}$. Far above $2^{64}$ two things change. The
sieving primes reach $3\cdot 10^{12}$ at $F = 10^{25}$ and $10^{16}$ at $F = 10^{32}$, so most
of them exceed the window and hit it at most once; and the phase $F \bmod p$ of each prime
needs 128-bit arithmetic. A sieve that visits every prime then pays per prime, not per hit.

The PrimeAtlas engines produce prime data for windows at such floors. Their \emph{jump
gear} (engines v7.0--v7.4) asks the opposite question: which large $q$ have a multiple in the
window. It answers it from the window's position, block by block, without visiting the
others. Helfgott~\cite{helfgott} gave this method first, with continued fractions in place
of the three-gap steps used here, and with an implementation tested up to $5\cdot 10^{18}$.
Section~\ref{sec:prior} compares the two.

\subsection*{Contributions}
\begin{enumerate}
  \item[(C1)] A Lean proof that the C code of the gear, transcribed with explicit
    wraparound and a double-precision rounding model, finds every landing of a block with
    its exact offset, under stated preconditions (Theorem~\ref{thm:main}).
  \item[(C2)] The engine bound $30F < 2^{128}$ that the proof requires, and a kernel-checked
    instance showing that the earlier guard $F < 2^{126}$ was too weak
    (Theorem~\ref{thm:overflow}). The engines were changed to the proved bound.
  \item[(C3)] A certificate for the three-gap steps that holds regardless of how the steps
    were obtained (Theorem~\ref{thm:cert}), used by engine v7.4 to skip most block setups.
  \item[(C4)] Optimality of the engines' gear threshold and block factor within their cost
    model (Theorem~\ref{thm:cost}).
\end{enumerate}
Measurements of the engines are reported in Section~\ref{sec:bench}; they are not part of
the formal results.

\section{Setting}\label{sec:setting}

Throughout, $F \ge 1$ is the lower end of the window, $W \ge 1$ its width, $M = 2^{64}$, and
$x \bmod m$ is the remainder of $x$ on division by $m$.

\begin{definition}[Landing]
For $q \ge 1$ let $g(q) = (q - F \bmod q) \bmod q$. The number $q$ \emph{lands} if
$g(q) < W$; then $F + g(q)$ is the least multiple of $q$ that is $\ge F$, and it lies in the
window.
\end{definition}

The gear walks blocks $q = Q + 30j$, $0 \le j < J$, in each of the eight residue classes
coprime to $30$, composite $q$ included. Its fixed-point quantities are
\[
  X_0 = \Big\lfloor \frac{F M}{Q} \Big\rfloor,\qquad
  S = \Big\lfloor \frac{a}{Q} \Big\rfloor M + \Big\lfloor \frac{(a \bmod Q) M}{Q}\Big\rfloor
    + \Big\lfloor \frac{r_1 M}{Q^2}\Big\rfloor,\quad a = \Big\lfloor \frac{30F}{Q}\Big\rfloor,\
    r_1 = 30F - aQ,
\]
so that $0 \le FM/Q - X_0 < 1$ and $0 \le 30FM/Q^2 - S < 2$. With $\fn{slack} = 2J + 8$,
\[
  A = (\fn{slack} - X_0) \bmod M,\qquad B = S \bmod M,\qquad v(j) = (A + Bj) \bmod M .
\]

\begin{definition}[Candidate condition and threshold]
The index $j$ is a \emph{candidate} if $v(j) < T$. The engines compute
\[
  t_d = \frac{W}{Q}\,M + 1.001\cdot\frac{F\,d_{\max}^2}{Q^3}\,M + 2\,\fn{slack} + 1,
  \qquad d_{\max} = 30(J-1),
\]
in double precision and set $T = \lfloor t_d \rfloor$; when $t_d \ge 1.8\cdot 10^{19}$ every
$j$ of the block is checked directly.
\end{definition}

\begin{definition}[Search and steps]
$\fn{first}(a, m, l, r)$ is the recursion
\[
\fn{first}(a, m, l, r) = \begin{cases}
 0 & l = 0,\\
 \text{none} & a \bmod m = 0,\\
 \lceil l/a \rceil & a\lceil l/a\rceil \le r,\\
 \big\lceil (l + m\,y)/a \big\rceil,\ y = \fn{first}(m \bmod a,\ a,\ (-r) \bmod a,\
   (-l) \bmod a) & \text{otherwise,}
\end{cases}
\]
with $a$ reduced mod $m$. For $T \le M$ let $g_a$ be the least $g \ge 1$ with
$Bg \bmod M < T$, $g_c$ the least $g \ge 1$ with $Bg \bmod M > M - T$,
$\alpha = Bg_a \bmod M$ and $\beta = M - (Bg_c \bmod M)$. The \emph{step rule} moves from a
candidate $v$ by $g_a$ if $v + \alpha < T$, by $g_c$ if $v \ge \beta$, and by $g_a + g_c$
otherwise.
\end{definition}

The engines find the first candidate $j_0$ with $\fn{first}$ on the arc shifted by $A$ (split
when it wraps), and the steps $(g_a, g_c)$ either by two more searches (v7.0, v7.1) or by one
subtractive Euclid walk on $d_a = Bg_a \bmod M$ and $d_c = (-Bg_c) \bmod M$ (v7.2), falling
back to the searches when the walk cannot decide. Each candidate is checked exactly: an
estimate $k = \lfloor (X_0 - Sj)/M \rfloor + 1$ of $\lceil F/q \rceil$ is corrected by two
loops to the least multiple of $q$ that is $\ge F$, and its offset is kept if it is below $W$.

\section{Results}\label{sec:results}

All C routines named below are in \lean{prime_sieve/engine_jump_gear.h} of the PrimeAtlas
repository~\cite{atlas}. Each is transcribed as a function on natural numbers in which every
128-bit operation is reduced mod $2^{128}$ and every 64-bit one mod $2^{64}$. The
double-precision threshold is modelled by the standard model of floating-point
arithmetic~\cite{higham}: every operation and every conversion from \lean{uint64_t} returns
$x(1+\delta)$ with $|\delta| \le 2^{-53}$, which holds for IEEE-754
round-to-nearest~\cite{ieee754} without underflow.

\subsection*{T1 The sieve argument}
\begin{theorem}\label{thm:sieve}
Let $S$ be a set of numbers that contains every prime $p$ with $p^2 < F + W$ that has a
multiple in $[F, F+W)$. Then a number $n \ge 2$ of the window is prime if and only if no
$d \in S$ with $2 \le d < n$ divides $n$. Moreover, for $q \ge W$ at most one multiple of $q$
lies in the window, it exists if and only if $g(q) < W$, and a multiple $m \ge F$ of any $q$
with $2 \le q < F$ is not prime.
\end{theorem}
\leanref{\lean{sieve_correct}, \lean{at_most_one}, \lean{landing_iff},
\lean{composite_harmless}}

So sieving by composite $q$ in addition to the primes does not change the output.

\subsection*{T2 The block routine}
\begin{theorem}\label{thm:main}
Assume
\[
 8 \le Q,\quad Q + 30J < 2^{64},\quad Q \le F < 2^{62} Q,\quad 60J \le Q,\quad
 1 \le W \le 2^{48},\quad 30F < 2^{128},
\]
let $t_d$ be computed as in the C code under the standard rounding model with
$t_d < 1.8\cdot 10^{19}$, $T = \lfloor t_d \rfloor$, and let $j_0$ and $(g_a, g_c)$ be the
results of the engine's searches. Then for every $j < J$ with $g(Q + 30j) < W$, the C loop
started at $(j_0, v(j_0))$ reaches $j$ before it leaves the block, the 128-bit estimate it
uses does not wrap, and the exact check returns $g(Q + 30j)$.
\end{theorem}
\leanref{\lean{jump_block_quick_exact} (v7.2: \lean{jump_first_range64} and
\lean{jump_gaps}), \lean{jump_block_exact} (v7.0, v7.1: three calls of
\lean{jump_first_range})}

The proof combines the following lemmas.

\begin{lemma}\label{lem:first}
For $l \le r < m$, $\fn{first}(a, m, l, r)$ is the least $x$ with $l \le ax \bmod m \le r$,
and it is none exactly when no such $x$ exists. The same holds for the shifted, possibly
wrapped arc, and for the 64-bit variant of v7.2 with results $\ge 2^{64}$ dropped.
\end{lemma}
\leanref{\lean{firstIn_spec}, \lean{firstRange_spec}, \lean{firstRange64_spec}}

\begin{lemma}\label{lem:step}
Let $T \le M$. From a candidate $v < T$ the step rule lands below $T$ and skips no $g$ with
$(v + Bg) \bmod M < T$. Hence, from the least candidate $j_0$, the steps visit every
candidate.
\end{lemma}
\leanref{\lean{step_lands}, \lean{step_first}, \lean{enumeration_complete}}

The lemma holds for every $T \le M$; the half-turn restriction often stated for the
three-gap step rule is not needed.

\begin{lemma}\label{lem:bridge}
If $g(Q+30j) < W$, $2j \le \fn{slack}$, $T \le M$ and
$T \ge \fn{slack} + 1 + E + (W-1)M/Q$ with $E \ge F(30j)^2M/Q^3$, then $v(j) < T$. Under the
standard rounding model, $1 \le W \le 2^{48}$, $Q \ge 1$ and $\fn{slack} \ge 3$, the computed
$t_d$ satisfies $t_d \ge \fn{slack} + 2 + F d_{\max}^2 M/Q^3 + (W-1)M/Q$.
\end{lemma}
\leanref{\lean{landing_candidate_fixed}, \lean{t_double_ok}}

\begin{lemma}\label{lem:exact}
If $8 \le Q < 2^{64}$, $F < 2^{62}Q$ and $30F < 2^{128}$, the 128-bit expressions for $X_0$
and $S$ equal their exact values. The two correction loops of the exact check return
$g(q)$ for every starting estimate, and with $60j \le Q$ the estimate $X_0 - Sj$ does not
underflow and every intermediate $kq$ stays below $2^{128}$.
\end{lemma}
\leanref{\lean{cX0_eq}, \lean{cS_eq}, \lean{cLanding_eq}, \lean{estimate_safe}}

\begin{lemma}\label{lem:gaps}
When the subtractive walk of \lean{jump_gaps} succeeds, its result is $(g_a, g_c)$.
\end{lemma}
\leanref{\lean{jumpGaps_spec}}

The walk keeps $d_a \equiv Bq_a$, $d_c \equiv -Bq_c \pmod M$, $q_a d_c + q_c d_a = M$ and
$d_a + d_c \ge T$. The first three say that $(q_a, d_a)$ and $(q_c, -d_c)$ are a basis of the
lattice $\{(g, r) : r \equiv Bg \bmod M\}$ of determinant $M$; a smaller step into $[0, T)$
would give a lattice point excluded by the sign cases. The argument uses only the
invariants, which gives the next result.

\subsection*{T3 The engine bound}
\begin{theorem}\label{thm:overflow}
For $F = 1.2\cdot 10^{37}$ and $Q = 3\cdot 10^{18} + 1$ one has $F < 2^{126}$,
$\lfloor F/2^{62}\rfloor + 1 \le Q$ and $Q^2 \le F$, and the C value of $S$ differs from the
exact one.
\end{theorem}
\leanref{\lean{cS_overflow} (by \lean{decide})}

At this $F$ the earlier guard $F < 2^{126}$ enabled the gear for this $Q$, but
$30F \ge 2^{128}$, the product \lean{30 * first} wraps, and landings are lost. With the guard
$F \le \lfloor (2^{128}-1)/30 \rfloor$ the engines' range and that of
Theorem~\ref{thm:main} coincide.

\subsection*{T4 A certificate for the steps}
\begin{theorem}\label{thm:cert}
Let $g_a, g_c \ge 1$, $d_a = Bg_a \bmod M$, $d_c = M - (Bg_c \bmod M)$. If
\[
 1 \le d_a < T,\qquad d_c < T,\qquad T \le d_a + d_c,\qquad g_a d_c + g_c d_a = M,
\]
then $g_a$ and $g_c$ are the least three-gap steps for $(B, T)$.
\end{theorem}
\leanref{\lean{gaps_certificate}}

\subsection*{T5 The cost model}
The engines split the $q$ between the classic path, costing about $1/\ln q$ per unit of $q$,
and the gear, costing about $C/q$ with $C = a\sqrt{F/W} + bW$; in v7.2, at block factor $f$,
$C(f) = \alpha/f + \beta(1 + f^2/2)$ with $\alpha = 0.5\,a\sqrt{F/W}$ and
$\beta = bW/1.125$.

\begin{theorem}\label{thm:cost}
Let $p^* \ge e$ satisfy $p^* = C\ln p^*$.
\begin{enumerate}
  \item For $q \ge e$ the gear is the cheaper path exactly for $q \ge p^*$, so the split at
    $p^*$ is the cheapest of all assignments of a finite set of $q$ to the two paths.
  \item For $C \ge 3$ a fixed point with $C \le p^* \le C^3$ exists, and the iteration
    $p \leftarrow \max(C\ln p, 3)$ from $\max(C, 3)$ rises to it with error at most
    $(1/\ln C)^k p^*$ after $k+1$ steps.
  \item $f^* = (\alpha/\beta)^{1/3}$ minimises $C(f)$ over $f > 0$, also after clamping to
    $[0.05, 64]$, and $f^*$ with the threshold of $C(f^*)$ costs no more than any block factor
    with any assignment.
\end{enumerate}
\end{theorem}
\leanref{\lean{gear_cheaper_above}, \lean{classic_cheaper_below}, \lean{threshold_optimal},
\lean{fixedPoint_exists}, \lean{iterP_error}, \lean{blockCost_min},
\lean{blockCost_clamp_min}, \lean{model_optimal}}

The constants $a$ and $b$ are fitted to measurements (defaults $a = 300$, $b = 0.1$); that
the model describes the machine is not proved. Neither the model nor its constants affect
correctness, since Theorem~\ref{thm:main} holds for every threshold.

\section{Use in the engines}\label{sec:engines}

The formal results changed the engines in two commits of the PrimeAtlas
repository~\cite{atlas}:
\begin{itemize}
  \item \lean{27fe0f5} (engines v7.0--v7.2): the gear is enabled only for
    $F \le \lfloor (2^{128}-1)/30 \rfloor$ and $W \le 2^{48}$, the hypotheses of
    Theorem~\ref{thm:main}; above them the classic path sieves. A block with
    $F \ge 2^{62}Q$, which the earlier code left without landings, now checks every $q$, and
    the block length is capped at $J \le Q/60$. Identity tests against the classic engine are
    unchanged bit for bit.
  \item \lean{ab4491b} (engine v7.4): each thread keeps the previous block's $(g_a, g_c)$ and
    reuses it when the conditions of Theorem~\ref{thm:cert} hold, checked with exact 128-bit
    products; otherwise it runs the walk. The C check transcribes the hypotheses of the
    theorem directly; that transcription is not part of the formalization. The pair is kept
    for about 57\,\% of blocks at v7.2's block factors. At $10^{11}$ numbers per run v7.4 is
    5--13\,\% faster than v7.2 on floors 27--30 and equal to it within the spread below
    (Section~\ref{sec:bench}).
\end{itemize}

\section{Relation to prior work}\label{sec:prior}

Helfgott~\cite{helfgott} sieves an interval of half-width $\Delta$ around $n$ classically up to
$K\Delta$ with $K \ge 5/2$, and above that replaces $n/m$ by its tangent line on blocks of
length $R = \lfloor M\sqrt{\Delta/4n}\rfloor$, sieving by all integers $m$ in that range. He
finds the candidates of a block by continued fractions, computing the modular inverse with
them, and reaches time $O(\Delta \log n)$ when $\Delta \ge n^{1/3}(\log n)^{2/3}$. His
implementation is a proof of concept tested up to $n = 5\cdot 10^{18}$; he remarks that above
$2^{64}$ highly optimized implementations of the traditional algorithm are harder to find for
comparison. The mathematical core of the jump gear is his method.

The PrimeAtlas engines differ in the following points. They run the method from $10^{18}$ to
$10^{30}$ in 64.64 fixed point with 128-bit arithmetic; they enumerate candidates by
three-gap steps rather than continued fractions; they combine the gear with a bucketed
classic path in the manner of Oliveira e Silva, Herzog and Pardi~\cite{oes}; and the split
point is not a multiple of the window but the fixed point of the cost model of
Theorem~\ref{thm:cost}, whose block term $a\sqrt{F/W}$ dominates for narrow windows. The runs
of $10^9$ numbers at $F = 10^{26}$ in Section~\ref{sec:bench} lie below Helfgott's condition
($n^{1/3}(\log n)^{2/3} \approx 7\cdot 10^9$ there).

The three-gap theorem is due to S\'os~\cite{sos} and \'Swierczkowski~\cite{swierczkowski}; we
use Slater's form~\cite{slater} for the gaps between successive returns to an interval.
Galway~\cite{galway} reduces the space of the sieve of Atkin and Bernstein. For narrow
windows at deep floors another route to proven primes is Sorenson's pseudosquares prime
sieve~\cite{sorenson2006}.

\section{Origin and measurements}\label{sec:bench}

The measurements in this section are not part of the formal results. All runs used an AMD
Ryzen~9 5900X (12 cores, 24 threads), 64\,GB RAM, WSL2, with 24 threads for every program.

\subsection*{The profile that led to the gear}
Engine v6.1 sieves a wheel-30 bitmap with primes from primesieve~\cite{primesieve}, small
primes in cache-sized segments and large ones through buckets~\cite{oes}. At
$F = 10^{25}$, $W = 2\cdot 10^{11}$ it spent 2--3\,s on start-up, 11--14\,s in the segmented
stage and 120--130\,s in the bucket stage, about 90\,\% of the sieve. There are about
$1.1\cdot 10^{11}$ primes below $\sqrt{F+W}$; a prime $p > W$ hits the window with
probability $W/p$, but its generation and phase are paid either way. This led to asking
which large $q$ land instead of where each prime lands.

\subsection*{The threshold}
On runs of $10^{10}$ numbers (v7.0 with a manual threshold, before the cost model existed),
the best threshold grew by about $\sqrt{10}$ per floor, as the block term predicts:

\begin{center}
\begin{tabular}{rrrrr}
\toprule
floor & no gear (v6.1) & $p^* = 2.6\cdot10^{11}$ & best $p^*$ & best time \\
\midrule
26 & 217.6\,s & 116.4\,s & $8.2\cdot10^{11}$ & 101.5\,s \\
27 & 739.3\,s & 373.9\,s & $2.9\cdot10^{12}$ & 263.9\,s \\
28 & 2513.5\,s & 1178.3\,s & $\ge 8.9\cdot10^{12}$ & 756.1\,s \\
\bottomrule
\end{tabular}
\end{center}

\subsection*{Speed}
The speeds come from one series of 204 runs on 16 November 2026, recorded with its script in
\lean{benchmarks/2026-11-16_engines} of the PrimeAtlas repository~\cite{atlas}. A run sieves
$N$ consecutive numbers from $F = 10^k$ (floor $k$) in windows of $10^7$. Each point is the
mean of two runs; on every floor the runs of engines v7.2 and v7.4 alternate in the order
v7.2, v7.4, v7.4, v7.2. The comparison program is Walisch's Pseudosquares Prime
Sieve~\cite{pss} (version 1.2, commit \lean{fccb96e}), which proves its output, in count
mode at $10^9$ numbers per floor, the size of its own benchmarks. The times leave out the
phase in which the engines write their windows: it depends on the disk, and the
pseudosquares sieve in count mode writes nothing. Prime counts agree between all runs of a
point and between the programs (for example $24\,127\,085$ at floor 18 and $15\,514\,691$ at
floor 28).

\begin{center}
\begin{tabular}{rrrrrr}
\toprule
 & \multicolumn{1}{c}{PSS} & \multicolumn{3}{c}{engine v7.4} & \multicolumn{1}{c}{v7.2} \\
\cmidrule(lr){2-2}\cmidrule(lr){3-5}\cmidrule(lr){6-6}
floor & $N = 10^9$ & $N = 10^9$ & $N = 10^{10}$ & $N = 10^{11}$ & $N = 10^{11}$ \\
\midrule
18 & 203 & 2054 & 7310 & 7920 & 7989 \\
19 & 145 & 1753 & 6224 & 8282 & 7707 \\
20 & 43.9 & 1324 & 4858 & 7838 & 7558 \\
21 & 41.4 & 959 & 3263 & 6831 & 7032 \\
22 & 36.8 & 632 & 1783 & 5043 & 5391 \\
23 & 34.2 & 383 & 1172 & 3744 & 3639 \\
24 & 29.7 & 200 & 855 & 2377 & 2349 \\
25 & 33.0 & 95.1 & 534 & 1594 & 1600 \\
26 & 28.5 & 42.7 & 308 & 1120 & 1083 \\
27 & 25.2 & 19.0 & 156 & 746 & 707 \\
28 & 25.4 & 8.2 & 73.8 & 458 & 407 \\
29 & 22.9 & 3.5 & 33.3 & 248 & 219 \\
30 & 21.7 & 1.5 & 14.7 & 124 & 115 \\
\bottomrule
\end{tabular}

\smallskip
{\small Millions of numbers per second, without the write phase.}
\end{center}

At $10^{11}$ numbers per run, the size it is built for, engine v7.4 is faster than the
pseudosquares sieve on every floor from 18 to 30. On runs of $10^9$ numbers it is faster up
to floor 26 and with $10^{10}$ up to floor 29: larger runs amortise the per-run cost of the
sieving primes. The engine uses up to 1.9\,GB for $10^9$ numbers, 8.2\,GB for $10^{10}$ and
11.2\,GB for $10^{11}$; the pseudosquares sieve 0.8\,GB. Runs shorter than about a minute
spread by up to 16\,\% between the two runs of a point; the spread of every point is in the
repository.

\section{Scope}\label{sec:scope}

The proof concerns a model of the C code written in Lean, not the compiled binary; the
agreement of model and code rests on the transcription, on the compiler, and on the engines'
identity tests, which compare whole runs bit for bit with the classic engine. Floating point
enters only through the standard rounding model. Outside the block routine the formalization
covers the sieve argument (Theorem~\ref{thm:sieve}) but not the mapping of offsets to the
wheel-30 bitmap or the bucket application, which the identity tests cover.

\section{The formalization}\label{sec:formalization}

The formalization is a Lean~4~\cite{lean4} project built against Mathlib~\cite{mathlib}, both
at version \lean{v4.35.0}: \lean{GearProof/JumpGear.lean} (correctness, about 1900 lines)
and \lean{GearProof/CostModel.lean} (cost model, about 300 lines, independent of the first).
\lean{lake build} completes with no \lean{sorry}, \lean{admit} or \lean{native_decide}, and the
theorems depend only on the axioms \lean{propext}, \lean{Quot.sound} and
\lean{Classical.choice}. To build:
\begin{verbatim}
lake exe cache get
lake build
\end{verbatim}

\section{Use of AI}

The question that led to the gear (asking which large $q$ land in the window instead of where
each prime lands, from the profile of engine v6.1), the gear's threshold from a cost model,
and the choice of statements to prove are the author's. Generative AI (Claude Opus 5.5,
Anthropic, used through Claude Code) contributed in the following ways:
\begin{itemize}
  \item It chose the technique that answers the question (tangent-line approximation,
    three-gap steps), drawing on its general knowledge of the literature, which includes
    Helfgott's method; it read Helfgott's paper after engine v7.0 was built. The author has
    not read it.
  \item It wrote the engine code and the Lean statements and proofs under the author's
    direction. The author checked the statements against the engines and by measurement; the
    correctness of the proofs rests on the Lean kernel.
  \item It drafted this text and checked the identifiers of the references. The author set
    the structure and the stated contributions.
\end{itemize}

\begin{thebibliography}{99}

\bibitem{atlas}
A.~Flamandzki,
\emph{PrimeAtlas},
\url{https://github.com/Flaman55/primeAtlas}.

\bibitem{bayshudson}
C.~Bays and R.~H.~Hudson,
\emph{The segmented sieve of Eratosthenes and primes in arithmetic progressions to $10^{12}$},
BIT 17 (1977), 121--127,
\href{https://doi.org/10.1007/BF01932283}{doi:10.1007/BF01932283}.

\bibitem{galway}
W.~F.~Galway,
\emph{Dissecting a sieve to cut its need for space},
Algorithmic Number Theory (ANTS-IV), Lecture Notes in Computer Science 1838, 2000, 297--312,
\href{https://doi.org/10.1007/10722028_17}{doi:10.1007/10722028\_17}.

\bibitem{helfgott}
H.~A.~Helfgott,
\emph{An improved sieve of Eratosthenes},
Math. Comp. 89 (2020), 333--350,
\href{https://doi.org/10.1090/mcom/3438}{doi:10.1090/mcom/3438};
\href{https://arxiv.org/abs/1712.09130}{arXiv:1712.09130}.

\bibitem{higham}
N.~J.~Higham,
\emph{Accuracy and Stability of Numerical Algorithms}, 2nd ed., SIAM, 2002,
\href{https://doi.org/10.1137/1.9780898718027}{doi:10.1137/1.9780898718027}.

\bibitem{ieee754}
\emph{IEEE Standard for Floating-Point Arithmetic}, IEEE Std 754-2019,
\href{https://doi.org/10.1109/IEEESTD.2019.8766229}{doi:10.1109/IEEESTD.2019.8766229}.

\bibitem{lean4}
L.~de~Moura and S.~Ullrich,
\emph{The Lean 4 theorem prover and programming language},
Automated Deduction -- CADE 28, Lecture Notes in Computer Science 12699, 2021, 625--635,
\href{https://doi.org/10.1007/978-3-030-79876-5_37}{doi:10.1007/978-3-030-79876-5\_37}.

\bibitem{mathlib}
The mathlib Community,
\emph{The Lean mathematical library},
Proceedings of the 9th ACM SIGPLAN International Conference on Certified Programs and Proofs
(CPP 2020), 367--381,
\href{https://doi.org/10.1145/3372885.3373824}{doi:10.1145/3372885.3373824}.

\bibitem{oes}
T.~Oliveira e Silva, S.~Herzog and S.~Pardi,
\emph{Empirical verification of the even Goldbach conjecture and computation of prime gaps up
to $4\cdot 10^{18}$},
Math. Comp. 83 (2014), 2033--2060,
\href{https://doi.org/10.1090/S0025-5718-2013-02787-1}{doi:10.1090/S0025-5718-2013-02787-1}.

\bibitem{primesieve}
K.~Walisch,
\emph{primesieve},
\url{https://github.com/kimwalisch/primesieve}.

\bibitem{pss}
K.~Walisch,
\emph{Pseudosquares Prime Sieve},
\url{https://github.com/kimwalisch/Pseudosquares-Prime-Sieve}.

\bibitem{slater}
N.~B.~Slater,
\emph{Gaps and steps for the sequence $n\theta \bmod 1$},
Proc. Cambridge Philos. Soc. 63 (1967), 1115--1123,
\href{https://doi.org/10.1017/S0305004100042195}{doi:10.1017/S0305004100042195}.

\bibitem{sorenson2006}
J.~P.~Sorenson,
\emph{The pseudosquares prime sieve},
Algorithmic Number Theory (ANTS-VII), Lecture Notes in Computer Science 4076, 2006, 193--207,
\href{https://doi.org/10.1007/11792086_15}{doi:10.1007/11792086\_15}.

\bibitem{sos}
V.~T.~S\'os,
\emph{On the distribution mod 1 of the sequence $n\alpha$},
Ann. Univ. Sci. Budapest. E\"otv\"os Sect. Math. 1 (1958), 127--134.

\bibitem{swierczkowski}
S.~\'Swierczkowski,
\emph{On successive settings of an arc on the circumference of a circle},
Fund. Math. 46 (1959), 187--189,
\href{https://doi.org/10.4064/fm-46-2-187-189}{doi:10.4064/fm-46-2-187-189}.

\end{thebibliography}

\end{document}
